\documentclass{article}

% -------------------------------------------------
% Packages
% -------------------------------------------------

\usepackage{xcolor}
\usepackage{amsmath}
\usepackage{amssymb}
\usepackage{amsthm}
\usepackage{graphicx}
\usepackage{subcaption}
\usepackage{matlab-prettifier}
\usepackage{listings}
\usepackage{blindtext}
\usepackage{geometry}
\usepackage{float}
\usepackage{dirtytalk}
\usepackage{enumitem}
\usepackage{derivative}
\usepackage{biblatex}
\usepackage{svg}
\usepackage{hyperref}

% -------------------------------------------------
% Colors
% -------------------------------------------------

\definecolor{darkgreen}{rgb}{0.0, 0.5, 0.0}

% -------------------------------------------------
% Hyperlinks
% -------------------------------------------------

\hypersetup{
    colorlinks=true,
    linkcolor=blue,
    citecolor=blue,
    urlcolor=blue
}

% -------------------------------------------------
% Theorem environments
% -------------------------------------------------

% Numbered
\newtheorem{theorem}{Theorem}[section]
\newtheorem{lemma}[theorem]{Lemma}
\newtheorem{proposition}[theorem]{Proposition}

% Unnumbered
\newtheorem*{theorem*}{Theorem}
\newtheorem*{lemma*}{Lemma}
\newtheorem*{proposition*}{Proposition}

\theoremstyle{definition}
\newtheorem{definition}[theorem]{Definition}
\newtheorem{example}[theorem]{Example}

\theoremstyle{remark}
\newtheorem{remark}[theorem]{Remark}

% -------------------------------------------------
% Equation numbering
% -------------------------------------------------

% Reset equation numbering at each subsection
\makeatletter
\@addtoreset{equation}{subsection}
\makeatother

% If you want equations to be numbered by subsection,
% uncomment the following line:
% \renewcommand{\theequation}{\thesubsection.\arabic{equation}}

% -------------------------------------------------
% Python / MATLAB code formatting
% -------------------------------------------------

\lstset{
    language=Python,
    basicstyle=\ttfamily\small,
    keywordstyle=\color{blue},
    stringstyle=\color{darkgreen},
    commentstyle=\color{gray},
    showstringspaces=false,
    frame=single,
    numbers=left,
    numberstyle=\tiny\color{gray},
    breaklines=true
}

% -------------------------------------------------
% Page geometry
% -------------------------------------------------

\geometry{
    a4paper,
    total={160mm,250mm},
    left=25mm,
    top=30mm,
    footskip=10mm
}

% -------------------------------------------------
% Document settings
% -------------------------------------------------

\setlength{\parindent}{0pt}

% -------------------------------------------------
% Title
% -------------------------------------------------

\title{Monte Carlo Simulation Uge 2}
\author{Thomas Sejer Canham}
\date{\today}

% =================================================
% Document
% =================================================

\begin{document}

% -------------------------------------------------
% Title page
% -------------------------------------------------

\begin{titlepage}

    \maketitle
    \thispagestyle{empty}

\end{titlepage}

\section{The Inverse Transfrom Method}

We discussed the inverse transform method. This method is usefull for generating random variables from a given distribution. However, we are restricted to distributions with a known cumulative distribution function (CDF) that is invertible. The method is based on the fact that if $U$ is a uniform random variable on the interval $[0, 1]$, then $X = F^{-1}(U)$ has the distribution function $F$. Pratically then, we can simulate a random variable $X$ with distribution function $F$ by generating a uniform random variable $U$ and then computing $X = F^{-1}(U)$.

\section{Discrete Distributions}
Sampling from is a bit different and requires algorithms that act on the uniform random variable $U$ to generate a discrete random variable $X$. To read more about this, refer to section 2.2 in the lecture notes.

\section{Location-Scale Families}
We say that a family of distributions is a location-scale family if a random variable $X$ with distribution $f(x)$
can be transformed into a random variable $Z = \mu + \sigma X$, such that $f_Z(z) = f(x; \mu, \sigma) := \frac{1}{\sigma} f(\frac{x - \mu}{\sigma})$ also is a well-defined distribution.


\end{document}
